% abc473 C - Change Schools の手書きメモの再現と振り返り
% ビルド: tools/build_thinking.sh abc473/c --preview
\documentclass[a4paper]{article}
\usepackage[margin=1.2cm]{geometry}
\usepackage{handmemo}
\usepackage{hyperref}
\pagestyle{empty}

\begin{document}

% ---------------------------------------------------------------- 1 枚目
\begin{memopage}{20260925\_174830.jpg}
  % クラスを袋で描いた
  \hw(4.3,2.4){生徒 $i$}
  \hwbag(3.1,3.75)(0.65,0.85)
  \hwbag(5.5,3.9)(0.55,0.85)
  \hwbag(7.8,4.1)(0.5,0.6)
  \hw(5.0,5.05){\normalsize $A_i$ class}
  \hw(4.7,6.2){\LARGE K class}
  \hw(4.8,7.45){N 人}

  % 3 つ目の袋へ向けた矢印とその元の語
  \draw[ink,hand,line width=0.7pt,-{Stealth[length=2mm]}] (8.55,2.75) to[bend right=25] (7.7,3.75);
  \hwunread(8.5,2.2){矢印の元の語（「lidong」のように見える）}

  \hw(8.6,4.1){人数max\ :\ \ 0}
  \draw[ink,hand,line width=0.7pt,-{Stealth[length=2mm]}] (11.05,5.1) -- (11.05,4.6);
  \hw(10.6,5.55){\normalsize count}
  \hwline(9.2,4.7)(9.2,7.2)
  \hw(8.7,8.3){人数最大のクラス}

  % 高橋くんが入ると +1
  \hw(1.4,10.75){+1}
  \hwline(1.45,11.25)(2.3,11.25)
  \redpen(3.0,10.2){入ったクラスは $c_i+1$、ほかは最大 $M$ のまま\\→ 喜ぶ条件は $c_i+1 \geq M$、つまり $c_i \geq M-1$}

  % 並べて top, next ...
  \hw(2.8,12.25){class num sort}
  \hw(8.2,12.2){top}
  \hwline(8.25,12.65)(9.4,12.65)
  \hw(9.6,12.1){next\ ..}
  \hwline(11.2,12.2)(12.0,12.15)
  \hwline(13.55,12.0)(14.0,12.1)
  \redpen(0.6,13.1){AC 版の {\hwlfont list(set(a))} は並びが保証されない\\人数 8, 1, 1 で 2 を出す（正しくは 1）}

  \hw(7.1,14.1){top, top−1 の個数}
  \redbox(6.9,13.45)(12.4,14.75)
  \redpen(7.0,15.4){結論は合っている。top−1 は「値 max−1」のこと\\→ コードでは {\hwlfont b[-2]}（2 番目に大きい人数）と読んで WA}
  \redarrow(9.2,15.0)(9.6,14.5)
\end{memopage}

% ---------------------------------------------------------------- 振り返り
\clearpage
\section*{abc473 C --- Change Schools}

\begin{itemize}
  \item 問題: \url{https://atcoder.jp/contests/abc473/tasks/abc473_c}
  \item 提出（2026-08-29, Python）:
    \href{https://atcoder.jp/contests/abc473/submissions/78822723}{WA} →
    \href{https://atcoder.jp/contests/abc473/submissions/78825493}{AC}
  \item 手書き原本: \texttt{Box/Photo Backup/手書き/atcoder/abc473/c/}（\texttt{20260925\_174830.jpg} の 1 枚。撮影 2026-09-25。前の 1 ページが再現。赤は後から入れた振り返り）
\end{itemize}

\subsection*{メモで考えたこと}
メモでは、クラスを袋として描いて「K class / N 人」「人数 max のクラス」「count」を置き、
高橋くんが入るとそのクラスが $+1$ になることを書いた。そのうえで「class num sort $\to$ top, next ...」と並べ、
「top, top $-1$ の個数」と結論を書いた。この結論（人数が max か max $-1$ のクラスの数）は正しい。

\subsection*{つまずき}
\begin{itemize}
  \item \textbf{WA}: コードでは「top $-1$」を\textbf{「2 番目に大きい人数」}として数えていた（\verb|b[-2]|）。
    2 番目に大きい人数が max $-1$ より小さいとき、そのクラスに入っても max のクラスに負けるので数えてはいけない。
    メモの「top $-1$」は値としての max $-1$ なのに、コードで「順位としての 2 番目」に読み替えてしまった。
  \item \textbf{AC 版にも潜在的なバグがある}: \verb|b = list(set(a))| はソートされる保証がない。
    小さい整数ではたまたま昇順に並ぶが、\verb|list(set([1, 8]))| は \verb|[8, 1]| になる。
    実際に \texttt{10 3 / 1 1 1 1 1 1 1 1 2 3}（人数 8, 1, 1）を AC 版に入れると \textbf{2} を出す（正しくは 1）。
    テストケースにこの形が無かったので通っただけ。
\end{itemize}

\subsection*{正しい考え方}
各クラスの人数を $c_i$、$M = \max c_i$ とする。クラス $i$ に入ると $c_i + 1$ になり、ほかのクラスは最大 $M$ のまま。
喜ぶ条件は $c_i + 1 \geq M$、つまり $\boldsymbol{c_i \geq M - 1}$。これを満たすクラスを数えるだけでよい（ソートも set も要らない）。
たとえば人数 5, 4, 2, 5, 3（$M=5$）なら、赤の点線（$M-1$）以上のクラスを数える。

\begin{center}
\begin{tikzpicture}[x=1cm,y=0.32cm]
  % 例: 人数 5, 4, 2, 5, 3（M = 5）。M-1 = 4 以上のクラスを赤で塗る
  \foreach \v/\col [count=\i] in {5/redpen!25,4/redpen!25,2/gray!25,5/redpen!25,3/gray!25} {
    \fill[\col] (\i*1.3-0.45,0) rectangle (\i*1.3+0.45,\v);
    \draw (\i*1.3-0.45,0) rectangle (\i*1.3+0.45,\v);
    \node[below] at (\i*1.3,0) {\small $c_{\i}=\v$};
  }
  \draw[dashed] (0.5,5) -- (7.3,5) node[right] {\small $M = 5$};
  \draw[redpen,dashed,line width=0.9pt] (0.5,4) -- (7.3,4) node[right] {\small $M-1 = 4$（ここ以上なら喜ぶ）};
\end{tikzpicture}
\end{center}

\begin{verbatim}
n, k = map(int, input().split())
c = [0] * k
for x in map(int, input().split()):
    c[x - 1] += 1
m = max(c)
print(sum(1 for v in c if v >= m - 1))
\end{verbatim}

\subsection*{次に活かすこと}
\begin{itemize}
  \item メモの結論を\textbf{条件式（$c_i \geq M - 1$）}の形で書いてからコードにする。
    「top $-1$ の個数」のような言葉のままだと、コードで別の意味に化ける。
  \item 並び順が欲しいときは \verb|sorted(set(a))|。\verb|set| の順序に頼らない。
\end{itemize}

\end{document}
